\chapter{Related Work}

This chapter sets out the background the project builds on. It starts with how
seed quality is checked today, by hand, in a working seed bank. It then surveys
five existing systems that automate part of that job, and notes where each one
falls short of what we need. From there it moves to the computer-vision ideas
behind our pipeline (object detectors and attention classifiers), the
synthetic-data techniques that let us train a detector without labelling
thousands of seeds by hand, and a short review of where AI sits in agriculture
more broadly. The chapter closes with a comparison of prior work against what
Seed Bank delivers.

\section{Traditional methods of seed analysis}

Traditional seed grading relies on mechanical systems that filter seed streams
through grids, sieves, and gravity tables, separating cohorts by size, weight,
and density into discrete quality tiers. These installations categorise large
volumes of seed into grades, typically spanning premium, first, second, and
lower-tier classifications, without requiring a human to examine each kernel.

These systems grade by physical proxies size, weight, and ensity so a
cracked or discoloured kernel that falls within the acceptable mass range passes
through undetected. Beyond that accuracy ceiling, the installations carry steep
drawbacks: the capital outlay for industrial sorting equipment runs to tens of
thousands of dollars, placing it out of reach for smaller seed banks, independent
laboratories, and individual farming operations. Ongoing specialist maintenance
is required to hold calibration, reconfiguring for a new crop or a revised
grading standard demands physical intervention, and the throughput advantage is
partly offset by the labour needed to operate and oversee the lines. The result
is a technology that serves large commercial operations but is inaccessible to
most of the agricultural sector.

\section{Existing solutions}

To understand the current landscape we looked at five systems, ranging from
industrial machinery to open-source research code. Each automates some part of
seed inspection, but most of them target growth potential (germination) or
species identification rather than the physical quality checks that long-term
storage actually requires.

\subsection{LemnaTec SeedAIxpert}

The LemnaTec SeedAIxpert \cite{lemnatec2026seedaixpert} is the reference point
for large corporate seed labs and plant breeders. It is an all-in-one hardware
and software product built around an imaging cabinet with controlled lighting,
so that every photograph is taken under identical conditions and the AI software
can produce reproducible scores. It is designed for high throughput, processing
large quantities of seed and replacing subjective human grading with
standardised numbers.

Relevance to our project: the strength of the system is also its barrier, which
is accessibility. It depends on custom, costly hardware and careful calibration
that is impractical for a typical university seed bank or a smaller conservation
effort. Its closed design also means we cannot adapt its algorithms to the
specific physical defects we care about. Our aim is to reach useful accuracy
from ordinary photographs without a dedicated imaging rig.

\subsection{PCS Agri Track}

PCS Agri Track \cite{pcsagri2026track} is a commercial product aimed at
nurseries, built to modernise inventory management and germination tracking.
Unlike the industrial lab equipment it runs on mobile and web, which lowers the
barrier to entry. It moves nurseries away from manual counting by digitising
the tracking process and recording data over time, such as germination curves.

Relevance to our project: the main drawback is its reliance on cloud processing
and a stable internet connection, which is not always available in the field.
Users have also reported that the interface is hard to follow and that the
models generalise poorly, working for common nursery seeds but struggling with
the diverse or rare species a seed bank holds. It still needs manual
intervention for edge cases, which undercuts the goal of automation.

\subsection{Seedy}

Seedy \cite{seedy2026app} is a mobile app for hobbyist gardeners, botanists, and
educators. Its purpose is on-demand seed identification and personal collection
management. A user photographs a seed with a phone camera, and the app's
recognition model returns the species along with metadata and growing guides.

Relevance to our project: Seedy shows that mobile seed analysis is both viable
and easy to use, which supports our own mobile capture flow. Its scope, though,
is strictly identification. It tells the user what a seed is, not how healthy it
is. It does not measure physical metrics or scan for damage and disease, which
is exactly the gap we are trying to fill.

\subsection{GerminationPrediction}

In the academic sphere the GerminationPrediction project of Grimm et al.
\cite{grimm2024germination} is a useful reference. It is open source, uses a
Faster R-CNN detector trained on a public dataset of over twenty-three thousand
images, and automates germination assessment by counting how many seeds in a
batch have sprouted.

Relevance to our project: this is the closest match technically, since it also
uses a CNN-based detector. The biological focus is different, though. They
measure viability (whether a seed will grow), while we measure physical
integrity (whether a seed is sound enough to store). A seed can be viable enough
to sprout and still carry cracks or fungal spores that make it risky to store in
a shared freezer. We borrow their open-source, reproducible approach but retrain
the focus toward defect detection rather than sprout counting.

\subsection{Multi-View Oil Palm system}

The Multi-View Oil Palm project \cite{oilpalmmvcnn} pushes computer vision past
the limits of a single image. It uses Multi-View CNNs (MVCNN) to analyse
irregularly shaped oil palm seeds from several angles, which gives it a 3D
understanding and strong accuracy on that crop.

Relevance to our project: this work makes a point we take seriously, that a
single top-down photograph can miss defects hidden on the side or underside of a
seed. Their multi-view setup gives richer data, but it is specialised for oil
palm and needs a complex rig to capture multiple angles. We aim for a middle
ground: high accuracy from accessible single-view photographs, without the
computational and hardware cost of a multi-view system.


\section{AI in agriculture}

AI is now a routine part of modern agriculture rather than an experiment. The
shift is driven by pressure on the sector: any land is shrinking and climate
conditions are more volatile, so the older "yield at all costs" mindset has given
way to a precision-first one, where efficiency and resilience matter as much as
raw output. The survey by Kamilaris and Prenafeta-Bold\'u
\cite{kamilaris2018deepag} documents this move across agriculture, from yield
prediction to weed and crop classification to quality grading, all shifting from
hand-crafted features to learned representations with steady accuracy gains where
labelled data was available.

In precision farming, models combine satellite imagery, soil-sensor readings,
and local weather forecasts into field-level prescriptions for how much water or
input each area needs. Closer to our problem is fine-grained computer vision for
quality control. Where ordinary detection just identifies a crop type,
fine-grained methods pick up subtle within-class differences such as hairline
fractures, surface discoloration, or early-stage fungal damage that often escape
the human eye. Mohanty et al. \cite{mohanty2016plant} show both the promise and
the catch: their CNNs reached high accuracy at recognising crop diseases from
leaf photographs on a curated dataset, but accuracy dropped sharply on images
taken under conditions different from the training set.

That last result points to the barriers the field still faces. The first is
capital cost: sensors and autonomous machinery carry a high upfront price that
keeps smaller operations out. The second is the digital divide, a widening gap
between fully integrated commercial farms and smallholders. The third is the one
that shaped our own work most: dataset variability and domain generalization.
Seed diversity and lighting changes break models trained on a narrow set, and
the usual fix is to collect local data and retrain for the specific harvest. The
gap between a clean benchmark and field photographs is the central practical
problem for any agricultural vision system, and it is the reason we treat data
variety as a first-class concern rather than an afterthought.

\section{Comparative analysis}

The components above are individually well established. What distinguishes this
project is not a new model but how the pieces are assembled into one workable
pipeline, and the decision to remove the data-labelling bottleneck with a
purpose-built generator rather than working around it.

Most prior work in seed and grain vision is a single artifact: one classifier or
one detector, trained against a hand-labelled dataset and evaluated once. Such a
system answers one question, either "where are the seeds" or "is this seed good",
but not the combined question the domain actually asks, which is how many seeds
are in this photo and how many of them are bad. It is usually trained on a
hand-labelled dataset, which caps the dataset size and therefore the accuracy.

Seed Bank differs on each point. It chains a detector and a per-type classifier
into one detect-then-classify pipeline, so a single photograph yields a located,
per-seed good/bad verdict and an aggregate report. It supports multiple crop
types, each with its own classifier, with detections grouped by type before the
classification stage. And the labelling bottleneck is attacked at the source by
MultiSeedGen, which generates perfectly labelled multi-seed data with controlled
variety, so the detector can be trained at a scale that hand-labelling could not
reach. Table~\ref{tab:ch2-compare} summarises the contrast.

\begin{singlespace}
\begin{longtable}{p{0.30\textwidth} p{0.31\textwidth} p{0.31\textwidth}}
\caption{Typical prior work in seed and grain vision versus the Seed Bank platform.}
\label{tab:ch2-compare} \\
\toprule
\textbf{Dimension} & \textbf{Typical prior work} & \textbf{Seed Bank} \\
\midrule
\endfirsthead
\toprule
\textbf{Dimension} & \textbf{Typical prior work} & \textbf{Seed Bank} \\
\midrule
\endhead
\bottomrule
\endfoot
Scope &
Single model: a detector \emph{or} a classifier &
Two-stage detect-then-classify pipeline producing per-seed verdicts and an aggregate report \\
Quality question answered &
``where are the seeds'' \emph{or} ``is this seed good'' &
``how many seeds, and how many are bad'' over a mixed photo \\
Crop types &
Usually one species, fixed at training &
Multiple crop types, each with its own classifier; detections grouped by type \\
Training data &
Hand-labelled, which caps dataset size &
Synthetic cut-and-paste generation with free, exact labels and controlled variety mixed with real data \\
Data integrity &
Leakage and label errors are easy to miss &
Leakage-safe split at the source-seed level, with controlled overlap and box clipping \\
\end{longtable}
\end{singlespace}

The wider point is that the contribution sits at the system level. Each model
family here is taken from published work and used as intended. The originality is
in binding detection, per-type classification, and a synthetic-data generator
into one tool that a non-researcher can run, and in attacking the labelling
bottleneck head-on rather than living with a small hand-labelled dataset. The
chapters that follow describe how that system is built and how it performs; the
experimental numbers are reported in the evaluation in Chapter~\ref{ch:testing}.
